

\centerline{\Huge Условия и решения задач.}
\centerline{\bf \Huge 7 класс}

\problem{7.1} Папа три года назад был в 9 раз старше дочки. Сейчас папа в 5 раз старше дочки. Сколько лет маме, если сумма возрастов мамы, папы и дочки сейчас равна 62 года. \par

\textbf{Решение.} Обозначим за $x$ возраст дочки 3 года назад. Тогда папе 3 года назад было $9x$. Стало быть сейчас ему $9x+3$, а дочке $x+3$. Сейчас папа в 5 раз старше дочки. Составим уравнение $9x+3 = 5(x+3)$, откуда получаем, что $x = 3$. Значит дочке сейчас 6 лет, а папе 30. Значит маме $62-30-6=26$ лет. \\ 

\textbf{Критерии.} \\
\textit{7 баллов} $-$ Полностью верное решение. \\
\textit{4 балла} $-$ Правильно составлено уравнение, связыващее возраст папы и дочки. \\
\textit{0-2 балла} $-$ Есть какие-то продвижения, которые могут привести к решению. \\ 
\textit{0 баллов} $-$ Решение отсутствует. Решение неверное, продвижения отсутствуют. \\ 

\problem{7.2}  Два программиста написали игру на языке Python. Правила игры следующие:\par
На мониторе написано число 2023. Каждый игрок по очереди может отнимать ненулевую цифру из числа находящегося на мониторе. Побеждает тот, после чьего хода на мониторе окажется число 0. У кого из игроков есть выигрышная стратегия, вне зависимости от ходов противника? У того, кто ходит первым или вторым?  \par

\textbf{Решение.} Будем называть число \textit{круглым}, если оно заканчивается на 0. \par
Приведём выигрышную стратегию для первого игрока: путём вычитания последней цифры он всегда будет получать из записанного на доске числа круглое число. \par
Своим первым ходом он вычитает из 2023 цифру 3 и получает круглое число 2023. Далее каждым своим ходом второй игрок будет вычитать ненулевую цифру из круглого числа, и будет получать не круглое число. А следующим своим ходом первый игрок легко будет получать круглое число, вычитая из написанного перед ним числа его последнюю ненулевую цифру.\\
Ясно, что рано или поздно игра закончится, и второй игрок не сможет выиграть, получив ноль, т.к. после каждого своего хода будет получать не круглое число. Значит, выиграет первый игрок.\\
\\ 

\textbf{Критерии.} \\
\textit{7 баллов} $-$ Полностью верное решение. \\
\textit{5 баллов} $-$ Приведена верная стратегия для первого игрока, но не доказано, почему он при ней выигрывает. \\ 
\textit{0-2 балла} $-$ Есть какие-то продвижения, которые могут привести к решению. \\ 
\textit{0 баллов} $-$ Решение отсутствует. Решение неверное, продвижения отсутствуют. \\ 

\problem{7.3} Лера, Аркадий и Владимир записали на доске по одному числу $-$ 0, 2 и 5 соответственно. Спустя минуту каждый из них одновременно стирает свое число и записывает вместо него сумму двух других(то есть Лера запишет 7, Аркадий 5, а Владимир 2). Спустя ещё минуту они повторяют это действие и так несколько раз. Могла ли сумма чисел Леры, Аркадия и Владимира в какой-то момент стать равной 20232022? \par

\textbf{Решение.} Давайте будем следить за суммой всех трех чисел. Если изначально были числа $a, \ b$ и $c$, то через минуту вместо них будут $b+c, \ a+c,$ и $a+b$ соответсвенно. В первый раз сумма трех чисел была равна $a+b+c$, а через минуту уже $(b+c) + (a+c) + (a+b) = 2(a+b+c)$. То есть каждую минуту сумма всех трёх чисел удваивается. Тогда через $n$ минут сумма будет равна $2^n(0+2+5) = 7\cdot2^n$. Но такое число не может быть равно $20232022$ по разным причинам. Например, 20232022 не делится на 7 или не делится на 4.\\

\textbf{Критерии.} \\
\textit{7 баллов} $-$ Полностью верное решение. \\
\textit{5 баллов} $-$ Доказано, что сумма трех чисел увеличивается в два раза с каждым разом. \\
\textit{0-2 балла} $-$ Есть какие-то продвижения, которые могут привести к решению. \\ 
\textit{0 баллов} $-$ Решение отсутствует. Решение неверное, продвижения отсутствуют. \\ 

\problem{7.4} Докажите, что для любого натурального $n$ число $11^n + 1$ не делится на $202300$.\par

\textbf{Решение.} Докажем, что наше число не делится даже $8$. Будем использовать арифметику остатков. Число 11 дает такой же остаток как и 3 при делении на 8. Значит достаточно доказать для $3^n + 1$. Число $3^n$ дает остаток либо 3 (в нечетной степени) либо 1 (в четной степени). А значит $3^n + 1$ дает остаток либо 2 либо 4 при делении на 8. Что и требовалось доказать. \\ 

\textbf{Критерии.} \\
\textit{7 баллов} $-$ Полностью верное решение. \\
\textit{2 балла} $-$ Сказано, что число не будет делиться на 8, но не доказано почему. \\
\textit{0-2 балла} $-$ Есть какие-то продвижения, которые могут привести к решению. \\ 
\textit{0 баллов} $-$ Решение отсутствует. Решение неверное, продвижения отсутствуют. \\ 

\problem{7.5} Хитрый Кощей Бессмертный взял в плен в свою пещеру 21 богатыря и нарисовал на земле окружность длины 1. После чего дал богатырям испытание: «Если вы сможете встать на края окружности так, что я не найду среди вас хотя бы 100 пар богатырей, расстояние между которыми не больше $\dfrac{1}{3}$, то я вас отпущу». Докажите, что как бы богатыри не встали, им не удасться спастись от хитрого злодея. 
\par
\textit{Расстоянием} между двумя богатырями будем считать длину кратчайшей дуги от одного богатыря к другому.

\textbf{Решение.} По индукции докажем утверждение для  $2n + 1$  точки(это наши богатыри) и $n$ дуг(это пары богатырей, которые образуют дугу на окружности). \textit{База}  n = 0.
\par 
\textit{Шаг индукции.} Пусть это верно для $n$. Докажем для $n+1$. Пусть имеется  $2n + 3$  точки. Рассмотрим "длинную" (более $\dfrac{1}{3}$ от окружности) дугу $AB$ с концами в этих точках (если таких нет, то "коротких" дуг более чем достаточно). Пусть $С$ – произвольная из оставшихся точек. По крайней мере одна из дуг $AС$, $BС$ не превосходит $\dfrac{1}{3}$. Действительно, так как дуги $AB, \ BC$ и $AC$ в сумме дают всю окружность(то суммарно по длине равны 1), не может оказаться такого, что все три дуги будут больше $\dfrac{1}{3}$. Значит какая-то из них не больше $\dfrac{1}{3}$. Проделаем так с каждой из $2n+1$ точек и найдем $2n+1$ новых дуг не больше $\dfrac{1}{3}$. Суммарно вместе с предположением индукции имеем $n^2 + 2n + 1 = (n+1)^2$

\textbf{Критерии.} \\
\textit{7 баллов} $-$ Полностью верное решение. \\
\textit{6 баллов} $-$ Доказано, что в графе из n вершин без треугольников не более чем $\dfrac{n^{2}}{4}$ ребер. \\ 
\textit{3 балла} $-$ Попытка доказательства общего случая по индукции, но с ошибкой. \\ 
\textit{2 балла} $-$ Приведен пример, когда таких расстояний ровно 100.\newpage
\textit{2 балла} $-$ Доказано аналогичное утверждение для меньшего числа богатырей (не меньше 5). \\
\textit{0-2 балла} $-$ Есть какие-то продвижения, которые могут привести к решению. \\ 
\textit{0 баллов} $-$ Решение отсутствует. Решение неверное, продвижения отсутствуют. \\ 
